\chapter{\IfLanguageName{dutch}{Long list}{Long list}}%
\label{ch:longlist}

Op basis van de requirementsanalyse in Hoofdstuk~\ref{ch:requirementsanalyse} worden in dit hoofdstuk alle mogelijke alternatieven per variabele component van de \gls{rag}-pipeline in kaart gebracht. Het doel van de long list is om via een literatuuronderzoek zo veel mogelijk kandidaten te identificeren. Per component worden alle gevonden alternatieven besproken, inclusief de opties die uiteindelijk in de short list zullen terechtkomen. Vervolgens worden ze samengebracht in een vergelijkingstabel die toont in welke mate elk alternatief voldoet aan de gedefinieerde requirements. Op basis van die tabel wordt per component duidelijk welke drie alternatieven het sterkst scoren en bijgevolg worden overgenomen in de short list.

\section{chunkingstrategieën}
\label{sec:longlist-chunking}

Chunking is de eerste en bepalende stap in de indexeringsfase van een \gls{rag}-pipeline. De keuze van chunkingstrategie heeft een directe impact op de kwaliteit van de retrieval: te kleine chunks verliezen de semantische samenhang van een passage, terwijl te grote chunks onnodig ruis introduceren in de context die aan het \gls{llm} wordt meegegeven~\autocite{Moro2025}. De literatuur identificeert de volgende families van chunkingstrategieën als mogelijke kandidaten.

\subsection{Fixed-size chunking}

Fixed-size chunking is de meest elementaire strategie waarbij tekst wordt opgesplitst in blokken van een vaste tokenlengte of aantal woorden, vaak aangevuld met een configureerbare overlap om de continuïteit te waarborgen. Het voordeel is de deterministische en snelle uitvoering zonder extra modelaanroepen, wat de reproduceerbaarheid maximaliseert~\autocite{QuEtAl2024}.~\textcite{BhatEtAl2025} onderzochten de optimale chunk-grootte systematisch over zes datasets en stelden vast dat 64 tot 128 tokens optimaal zijn voor feitelijke vraag-en-antwoordtaken, terwijl 512 tot 1.024 tokens beter scoren voor contextueel begrip. De voornaamste beperking is echter dat de strategie de tekststructuur volledig negeert, waardoor een chunk midden in een zin of kritieke tabel kan worden afgebroken, wat de semantische integriteit schaadt. Het vormt echter wel een goede baseline om andere chunkingstrategieën naast te vergelijken.

\subsection{Sentence-based chunking}

Bij deze methode wordt de tekst opgesplitst op zinsgrenzen. Hoewel dit garandeert dat elke chunk een grammaticaal volledige zin bevat, is de beperking voor complexe documenten aanzienlijk: veel zinnen missen de noodzakelijke context zonder hun omliggende alinea. Dit leidt tot een lage \textit{context recall}, waarbij de opgehaalde informatie onvoldoende contextuele samenhang bevat voor het \gls{llm} om een feitelijk antwoord op te baseren~\autocite{LatifEtAl2025}.

\subsection{Recursive character chunking}

Deze strategie splitst tekst hiërarchisch op basis van een geordende reeks separatoren (zoals dubbele en enkelvoudige regelafstanden) totdat de gewenste chunk-grootte bereikt is~\autocite{LangChainDocs2024}. Het respecteert hiermee de natuurlijke alinea- en lijststructuur van documenten zonder afhankelijk te zijn van extra computationele kosten.~\textcite{Amiri2025} evalueerden 25 chunking-configuraties over vijf methodefamilies en 48 embeddingmodellen in een domeinspecifieke context en stelden vast dat recursive token-based chunking consistent het best presteerde, met tot 45\% hogere domain-weighted precision dan de beste fixed-size variant. Ook~\textcite{SmithTroynikov2024} bevestigen in hun evaluatiestudie dat de keuze van chunking-strategie tot 9\% verschil in recall veroorzaakt, waarbij recursive character splitting in de praktijk consistent goed presteert. Hoewel het de industriestandaard is vanwege zijn betrouwbaarheid~\autocite{Amiri2025}, blijft de zwakte dat de methode puur op basis van tekens werkt en geen fundamenteel begrip heeft van de semantische overgangen binnen de tekst.

\subsection{Semantic chunking}

Semantic chunking, geformaliseerd door~\textcite{Kamradt2023}, berekent de cosinusgelijkenis - overeenkomst tussen 2 vectoren -  tussen opeenvolgende zin-embeddings en start een nieuwe chunk wanneer de betekenis verschuift~\autocite{Kamradt2023}. Dit resulteert in een zeer hoge interne coherentie en \textit{context precision} doordat elke chunk één onderwerp omvat~\autocite{MorenoCediel2026}. De keerzijde is de aanzienlijk hogere complexiteit dan vooraf besproken chunking strategieën en berekeningskost; voor elke zin in het corpus is een afzonderlijke embedding-aanroep vereist, wat het indexeringsproces vertraagt.

\subsection{Document structure chunking}

Deze aanpak gebruikt de opmaakstructuur — zoals koppen en paginagrenzen — als natuurlijke splitspunten. Hoewel dit theoretisch leidt tot zeer coherente segmenten, is de methode in de praktijk extreem fragiel. Bij de conversie van complexe \gls{esg}-rapporten naar platte tekst gaat opmaakinformatie vaak verloren, waardoor deze strategie te sterk afhankelijk is van een perfecte en foutgevoelige \gls{pdf}-parsing~\autocite{Amiri2025}.

\subsection{Vergelijkingstabel chunkingstrategieën}

Tabel~\ref{tab:longlist-chunking-tabel} brengt alle besproken chunkingstrategieën samen en toont in welke mate ze voldoen aan de technische requirements voor de PoC van Turtle Srl. De evaluatie van chunkingstrategieën focust op de balans tussen het behoud van context en de verwerkingssnelheid.

\begin{table}[ht]
\centering
\small
\begin{tabular}{lcc}
\toprule
\textbf{Strategie} & \textbf{CO (0-2)} & \textbf{SI (0-2)} \\
\midrule
Fixed-size chunking     & 0 & 0 \\
Recursive character     & 0 & 1 \\
Semantic chunking       & 1 & 2 \\
Sentence-based chunking & 0 & 1 \\
Document structure      & 0 & 1 \\
\bottomrule
\end{tabular}
\caption[Evaluatie van chunkingstrategieën op basis van overhead en semantiek.]{\label{tab:longlist-chunking-tabel}Vergelijking van chunkingstrategieën op een ordinale schaal (0-2).}
\vspace{2mm}
\footnotesize{\textbf{Legenda (CO):} 0 = Lineaire tekstverwerking zonder model-inferentie (\gls{cpu}-gebonden), 1 = Vereist inferentie van een embedding-model (GPU-acceleratie vereist), 2 = Vereist volledige \gls{llm}-inferentie per verwerkingseenheid (hoge latency en kosten). \\ \textbf{Legenda (SI):} 0 = Contextblinde splitsing op basis van vaste tekstlengtes, 1 = Behoud van syntactische of structurele grenzen (zinnen, paragrafen), 2 = Splitsing op basis van semantische coherentie of intentie-analyse.}
\end{table}

\textbf{Onderbouwing:}
\textit{Semantic chunking} wordt verkozen als het optimale evenwicht tussen intelligentie (SI=2) en schaalbaarheid (CO=1).

\section{Embeddingmodellen}
\label{sec:longlist-embedding}

Het embeddingmodel zet tekstchunks om in numerieke vectoren die in de vectordatabase worden opgeslagen. De kwaliteit van deze representaties bepaalt hoe goed de retriever semantisch verwante passages terugvindt bij een zoekopdracht. Mijn copromotor stelde specifieke eisen: modellen moeten meertalig zijn, geoptimaliseerd voor instruction-following en sterk presteren op retrieval op basis van zowel vragen als trefwoorden. Als objectief selectiecriterium wordt het \gls{mteb}-leaderboard gehanteerd, de academische standaard voor het vergelijken van embeddingmodellen over 58 datasets en 112 talen~\autocite{MuennighoffEtAl2023}.

\subsection{Snowflake arctic-embed-l-v2.0}

Snowflake arctic-embed-l-v2.0, beschreven door~\textcite{YuEtAl2024}, is een 303M-parameter embeddingmodel gebouwd op een XLM-R-Large backbone en uitgebracht onder de Apache 2.0-licentie. Het model werd specifiek ontworpen voor meertalige retrieval zonder kwaliteitsverlies op Engelstalige taken. De CLEF meertalige retrieval-benchmark, waarbij \gls{ndcg@10} gemeten wordt over Engels, Frans, Spaans, Italiaans en Duits, geeft een gemiddelde van 0,541 — significant hoger dan BGE-M3 (0,410) en multilingual-e5-large (0,431)~\autocite{YuEtAl2024}. Dit maakt het model de sterkste open-source keuze voor de Italiaans-Engelse documentbasis van Turtle Srl. De maximale inputlengte bedraagt 8.192 tokens en het model ondersteunt Matryoshka-stijl truncatie naar 256 dimensies met minder dan 3\% kwaliteitsverlies, wat de opslag in Qdrant aanzienlijk kan optimaliseren. Het integreert rechtstreeks via \texttt{SentenceTransformersDocumentEmbedder} in Haystack 2.x en kan lokaal worden ingezet op een Azure VM.

\subsection{BAAI BGE-M3}

BGE-M3, gepubliceerd door het Beijing Academy of Artificial Intelligence, introduceert een unieke combinatie van drie retrievalfunctionaliteiten via één modelgewicht: dense retrieval (semantisch zoeken via cosine-gelijkenis), sparse retrieval (trefwoordzoeken via geleerde gewichten) en multi-vector retrieval (ColBERT-stijl token-niveau matching)~\autocite{Chen2024}. Het is het eerste embeddingmodel dat deze drie functionaliteiten combineert. Verdere specificaties zijn ondersteuning voor meer dan 100 talen en een maximale inputlengte van 8.192 tokens. De multi-vector uitvoer sluit naadloos aan bij Qdrant's native ondersteuning voor hybride zoekopdrachten via \gls{rrf}~\autocite{QdrantHybrid2025}. Dit is bijzonder relevant voor \gls{esg}-vragenlijsten waarbij zowel semantisch als op trefwoord gezocht moet kunnen worden. De beperking is de infrastructuurkost: als open-source model vereist BGE-M3 lokale hosting of een containerized Azure-deployment.

\subsection{Microsoft multilingual-e5-large-instruct}

De E5-modelfamilie van Microsoft Research introduceerde het concept van \textit{query-} en \textit{passage-}prefixen als instruction-tuning mechanisme~\autocite{WangEtAl2022}. Multilingual-e5-large-instruct, beschreven door~\textcite{WangEtAl2024b}, initialiseert vanuit xlm-roberta-large, ondersteunt 100 talen en is getraind op 1 miljard meertalige tekstparen. Door simpelweg \texttt{query:} voor een vraag en \texttt{passage:} voor een documentchunk te plaatsen, dwingt men het model de asymmetrische relatie tussen vraag en antwoord te modelleren~\autocite{Wang2024a}. Vergeleken met BGE-M3 is dit model lichter in gewicht, wat snellere lokale hosting op een Azure-server mogelijk maakt. De beperking is de afwezigheid van hybride zoekfunctionaliteit.

\subsection{Cohere embed-multilingual-v3.0}

Cohere's embed-multilingual-v3.0 is een sterk commercieel model dat specifiek getraind is voor search- en retrieval-taken en uitstekend presteert op het \gls{mteb}-leaderboard voor meertalige taken~\autocite{MuennighoffEtAl2023}. Echter, de \gls{poc} maakt reeds gebruik van de Azure-infrastructuur met OpenAI-integratie. Het toevoegen van Cohere als derde externe \gls{api}-provider verhoogt de architectuurcomplexiteit zonder een unieke testdimensie toe te voegen ten opzichte van text-embedding-3-large, dat de proprietaire categorie reeds afdekt.

\subsection{Jina jina-embeddings-v2-base-multilingual} % TODO (v3 model !!!)

Het Jina-model ondersteunt contexten tot 8.192 tokens~\autocite{JinaAI2024}, wat interessant is voor lange \gls{esg}-documenten. BGE-M3 biedt echter dezelfde lange context én bovendien hybride zoeken, waardoor het Jina-model technologisch inferieur is voor de specifieke eisen van dit onderzoek, zonder compenserende voordelen.

\subsection{Nomic embed-text-v1.5}

Nomic embed-text-v1.5, mxbai-embed-large-v1 en GTE-large-en-v1.5 zijn Apache 2.0-modellen die sterk scoren op de Engelstalige \gls{mteb} Retrieval-benchmark. Zij worden echter uitgesloten omdat zij uitsluitend op Engelse documenten zijn getraind en dus niet voldoen aan de meertaligheidseis voor de Italiaans-Engelse documentbasis van Turtle Srl.

\subsection{Vergelijkingstabel embeddingmodellen}

De embeddingcomponent is verantwoordelijk voor de vectorisatie van de tekstfragmenten en bepaalt de effectiviteit van de zoekoperaties (retrieval) binnen de \gls{rag}-pipeline. Voor de evaluatie van deze modellen ligt de focus op de kwaliteit van de informatie-ophaling, de ondersteuning voor meertalige \gls{esg}-documentatie (Italiaans en Engels) en de infrastructurele conformiteit binnen de \gls{eu}-jurisdictie.

De kandidaten worden getoetst aan de volgende criteria:

\begin{itemize}
  \item \textbf{MTEB-Retrieval (nDCG@10):} De primaire metriek die meet hoe accuraat het model de meest relevante fragmenten bovenaan de resultaten plaatst.
  \item \textbf{Soevereiniteit Tier (ST):} Een numerieke schaal (0-3) voor juridische veiligheid:
    \begin{itemize}
      \item \textbf{0 (EU-Native):} Bedrijf en hosting volledig binnen de \gls{eu}.
      \item \textbf{1 (EU-Hosted):} US-origin model, maar door \textit{Open Weights} volledig te hosten op eigen \gls{eu}-servers.
      \item \textbf{2 (US-Cloud):} Propriëtaire \gls{api} (SaaS) onderhevig aan de Amerikaanse CLOUD Act.
      \item \textbf{3 (Niet-Westers):} Jurisdicties zonder \gls{gdpr}-adequaatheidsbesluit.
    \end{itemize}
  \item \textbf{Meertaligheid (ML):} Native ondersteuning voor zowel Italiaanse als Engelse tekstbronnen, essentieel voor de cross-linguale retrieval binnen de \gls{esg}-dataset.
  \item \textbf{Open-source (OS):} Beschikbaarheid van de modelparameters (\textit{Open Weights}) voor lokale implementatie, wat volledige controle over de dataverwerking mogelijk maakt.
  \item \textbf{Hybride Retrieval:} De mate waarin het model zowel semantische (dense) als lexicale (sparse) zoekopdrachten ondersteunt.
  \item \textbf{Context Window:} De maximale invoerlengte per tekstfragment in tokens.
\end{itemize}

Tabel~\ref{tab:embeddingmodels} brengt alle besproken embeddingmodellen samen en toont in welke mate ze voldoen aan de technische requirements voor de PoC van Turtle Srl.

\begin{table}[ht]
\centering
\small
\begin{tabular}{lcccccc}
\toprule
\textbf{Model} & \textbf{MTEB-R} & \textbf{ST (0-3)} & \textbf{ML (IT)} & \textbf{OS} & \textbf{Hybride} & \textbf{Context} \\
\midrule
BAAI BGE-M3                        & 54.60\% & 1 & Ja  & Ja  & Ja  & 8192 \\ % keuze over Jina -> hybrid verbeterd MTEB-R score (https://www.mdpi.com/2504-4990/8/4/91 | p15)
Snowflake Arctic-L-V2.0            & 58.36\% & 2 & Ja  & Ja  & Nee & 8192 \\ % keuse
MS Multilingual-E5-Large-Instruct  & 57.12\% & 2 & Ja  & Ja  & Nee & 512  \\ % keuze
Jina Embeddings-V3                 & 55.76\% & 0 & Ja  & Ja  & Nee & 8192 \\ % Anders geen hybrid
Nomic Embed-Text-V1.5              & 59.45\% & 2 & Nee & Ja  & Nee & 8192 \\ % Niet meertalig
Cohere Embed-Multi-V3              & 59.22\% & 2 & Ja  & Nee & Ja  & 512  \\ % niet open source
\bottomrule
\end{tabular}
\caption{Vergelijking van embeddingmodellen op basis van retrieval-nauwkeurigheid, meertaligheid en soevereiniteit.}
\label{tab:embeddingmodels}
\end{table}

\section{Large Language Models}
\label{sec:longlist-llm}

Het \gls{llm} vormt de generatieve kern van de \gls{rag}-pipeline. Het ontvangt de opgehaalde chunks als context en formuleert op basis daarvan een coherent antwoord. Binnen de \gls{rag}-architectuur fungeert het als een actieve verwerker van niet-parametrische data: in plaats van te vertrouwen op kennis die tijdens pre-training is opgeslagen, redeneert het model over de aangeleverde documentfragmenten~\autocite{Lewis2020}. De volgende kandidaten werden onderzocht.

\subsection{OpenAI GPT-4o}

% TODO

\subsection{Meta Llama 3.3 70B Instruct}

Llama 3.3 70B Instruct, gepubliceerd door Meta~\autocite{GrattafioriEtAl2024}, is een 70 miljard parameter model uitgebracht onder de Llama 3.3 Community License. Het behaalt een \gls{mmlu}-score van 86,0\%, MMLU-Pro van 68,9\% en een IFEval-score van 92,1\% — de hoogste instructieopvolgingsscore van alle onderzochte kandidaten, wat de betrouwbaarheid van prompt-instructies in de Haystack-pipeline maximaliseert. Op de Vectara Hallucination Leaderboard scoort het model 4,1\% hallucinaties~\autocite{TamberEtAl2025}, wat het tot het meest faithfulness-betrouwbare open-source model in de short list maakt. Het model ondersteunt officieel acht talen waaronder Italiaans en heeft een context window van 128.000 tokens. Llama 3.3 70B is beschikbaar op Microsoft Azure Foundry met EU-dataresidentie, wat Azure-compatibiliteit en \gls{gdpr}-naleving maximaliseert.

\subsection{DeepSeek DeepSeek-V3}

DeepSeek-V3 is een open-source model dat op benchmarks presteert op het niveau van GPT-4, tegen een fractie van de kosten. Dit model wordt uitdrukkelijk uitgesloten op grond van \gls{gdpr}-vereisten: DeepSeek is een Chinese aanbieder en het doorsturen van de vertrouwelijke \gls{esg}-documenten van Turtle Srl naar DeepSeek-servers is onverenigbaar met de Europese wetgeving inzake gegevensbescherming~\autocite{HuangEtAl2024}. Voor gevoelige bedrijfsdocumenten is datasoevereiniteit een harde uitsluitingsgrond.

\subsection{Google Gemma 3 27B IT}

Gemma 3 27B IT, uitgebracht door Google DeepMind~\autocite{GemmaTeam2025}, is een 27 miljard parameter instruction-tuned open-source model beschikbaar onder de Gemma-licentie, die commercieel gebruik toelaat. Het achtervoesel \textit{IT} staat expliciet voor \textit{Instruction Tuned}: het model is na pre-training gefinetuned op het opvolgen van complexe instructies, het genereren van gestructureerde output en vraag-en-antwoord-taken, wat naadloos aansluit bij de Metadata-Aware Prompt Builder in de Haystack-pipeline. Het model beschikt over een context window van 128.000 tokens en biedt meertalige ondersteuning, waaronder Italiaans en Engels~\autocite{GemmaTeam2025}. Het is beschikbaar via de Azure AI Foundry Model Catalog en kan volledig lokaal worden ingezet op een Azure VM via vLLM of Ollama, waardoor \gls{gdpr}-naleving gegarandeerd is zonder dat data de EU-infrastructuur verlaat. Met 27 miljard parameters biedt het model een gunstige latency-verhouding ten opzichte van grotere modellen in deze evaluatie.

\subsection{Mistral AI Mistral Small 4 (mistral-small-4-119b-2603)}

Mistral Small 4, uitgebracht in maart 2026 door Mistral AI~\autocite{MistralAISmall42026}, is een 119 miljard parameter instruction-tuned model beschikbaar onder de Apache 2.0-licentie. Het combineert de Europese datasoevereiniteitsgaranties van de Mistral-modelfamilie met een volwaardige instruction-tuned architectuur — het doorslaggevende onderscheid ten opzichte van Mistral Large 2, dat als basismodel de instruction-tuning-requirement mist. Als Frans bedrijf gevestigd in Parijs verwerkt Mistral AI alle data exclusief binnen de \gls{eu}, waardoor het niet valt onder de Amerikaanse CLOUD Act~\autocite{MistralAISmall42026,Rutherford2019}. Het model ondersteunt meerdere talen waaronder Italiaans en Engels, beschikt over een context window van 128.000 tokens en behaalt sterke prestaties op meertalige instructieopvolgings- en redeneer-benchmarks. De combinatie van EU-datasoevereiniteit, instruction-tuning, Apache 2.0-licentie en Azure-compatibiliteit via de Foundry Model Catalog en volledig self-hosted deployment op een Azure VM via vLLM plaatst dit model als de sterkste kandidaat voor de data-governance-dimensie van de vergelijking.

\subsection{Vergelijkingstabel LLMs}

De selectie van de generatieve component is gebaseerd op de balans tussen taalkundige intelligentie, het vermogen om strikte instructies op te volgen en de juridische datasoevereiniteit voor de verwerking van gevoelige bedrijfsinformatie van Turtle Srl.

De modellen worden beoordeeld op de volgende criteria:
\begin{itemize}
    \item \textbf{MMLU (Massive Multitask Language Understanding):} Een benchmark voor de algemene redeneervaardigheid en kennis van het model.
    \item \textbf{IFEval (Instruction Following Evaluation):} Een geverifieerde score die meet hoe accuraat een model strikte formele instructies opvolgt (essentieel voor JSON-outputs).
    \item \textbf{Soevereiniteit Tier (ST):} Dezelfde classificatie (0-3) als gehanteerd bij de embeddingmodellen, gericht op het voorkomen van ongeautoriseerde data-export.
    \item \textbf{Modeltype:} Het onderscheid tussen \textit{Open Weights} (lokaal hostbaar) en \textit{Proprietary} (SaaS-modellen).
\end{itemize}

Tabel~\ref{tab:llm-longlist} brengt alle besproken \glspl{llm} samen en toont in welke mate ze voldoen aan de technische requirements voor de PoC van Turtle Srl.

\begin{table}[ht]
\centering
\small
\begin{tabular}{lcccccc}
\toprule
\textbf{Model} & \textbf{MMLU} & \textbf{IFEval} & \textbf{Context} & \textbf{Type} & \textbf{ST} \\
\midrule
Mistral Small 4 (2603)  & 73.5\% & ?      & 32k  & Open Weights & 0 \\
Meta Llama 3.3 70B      & 86.0\% & 92.1\% & 128k & Open Weights & 1 \\
Google Gemma 3 27B IT   & 78.6\% & 90.4\% & 128k & Open Weights & 1 \\
GPT-4o (OpenAI)         & 85.7\% & 81.0\% & 128k & Proprietary  & 2 \\
DeepSeek-V3             & 88.5\% & 88.5\% & 128k & Open Weights & 3 \\
\bottomrule
\end{tabular}
\caption{Objectieve vergelijking van LLM-kandidaten op basis van intelligentie, stuurbaarheid en soevereiniteit.}
\label{tab:llm-longlist}
\end{table}